\chapter*{图片致谢}

本书中的插图为原创。照片依其各自的自由许可使用，谨向其作者致谢：

\begin{itemize}
  \item \emph{克洛德·贝尔纳}（第 1 章）：照片，19 世纪，维基共享资源，公有领域。
  \item \emph{软木细胞}（第 5 章）：罗伯特·胡克《显微图谱》（\emph{Micrographia}，1665）中的插图，维基共享资源，公有领域。
  \item \emph{安东尼·范·列文虎克}（第 5 章）：Jan Verkolje 绘制的油画，约 1680 年，经维基共享资源由荷兰国立博物馆提供，公有领域。
  \item \emph{Escherichia coli}（第 5 章）：扫描电子显微照片，NIAID，经维基共享资源提供，公有领域。
  \item \emph{莱纳斯·鲍林}（第 12 章）：照片，诺贝尔基金会，1962 年，经维基共享资源提供，公有领域。
  \item \emph{梅尔文·卡尔文}（第 14 章）：照片，劳伦斯伯克利实验室，经维基共享资源提供，公有领域。
  \item \emph{汉斯·克雷布斯}（第 15 章）：照片，诺贝尔基金会，1953 年，经维基共享资源提供，公有领域。
  \item \emph{人类核型}（第 17 章）：美国国家人类基因组研究所，经维基共享资源提供，公有领域。
  \item \emph{红细胞、血小板和淋巴细胞}（第 7 章）：扫描电子显微照片，美国国家癌症研究所电子显微镜设施，经维基共享资源提供，公有领域。
  \item 第~\ref{ch:b1:classifying-biodiversity}章：卡尔·林奈的肖像，Alexander Roslin 绘于 1775 年，斯德哥尔摩国家博物馆藏——公有领域（维基共享资源）。
  \item 第~\ref{ch:b1:classifying-biodiversity}章：卡尔·沃斯 的照片，摄影者 Don Hamerman——CC~BY~3.0（维基共享资源）。
  \item 第~\ref{ch:b1:classifying-biodiversity}章：恩斯特·海克尔《自然界的艺术形态》（\emph{Kunstformen der Natur}，1904）中的两幅图版，硅藻类（Diatomea）与笼形放射虫（Cyrtoidea）——公有领域（维基共享资源）。
\end{itemize}

每张照片的完整来源链接和许可证网址，均列于本书代码库中的 \texttt{images/book3/CREDITS.md}。

\bigskip
与照片和原创插图并列，全书各章中约七十幅图为 AI 生成插图，由 OpenAI 图像生成功能依本书编辑撰写的提示词制作，并在收录前经过生物学准确性审校。每幅生成图像的完整提示词均记录于代码库中的 \texttt{images/book3/ai/PROMPTS.md}。
\clearpage
